\subsection{Tuning the LBBNN baseline}
\label{sec:lbbnn-tuning}

Since LBBNN underperforms the sampling-based methods on the UCI
benchmarks, we verified that this reflects the method rather than a poor
choice of hyperparameters. We tuned the two knobs governing its
spike-and-slab posterior on \textsc{boston}, over the same five splits
used throughout, holding all else fixed.

\paragraph{Concrete temperature.}
The relaxation temperature $\tau$ of the Bernoulli inclusion variables was
varied over $\tau \in \{0.02, 0.05, 0.1, 0.2, 0.3, 0.5, 0.8\}$, against our
default $\tau = 0.2$. Within the range where the variational posterior
remains non-degenerate ($\tau \leq 0.2$), test RMSE spans $3.36$--$3.52$,
a band narrower than its own standard error ($\approx 0.23$) and not
monotone in $\tau$. The best setting, $\tau = 0.05$, gives
$3.36 \pm 0.24$ against $3.52 \pm 0.23$ at the default; the paired
difference over splits is $-0.16$ ($t = -1.6$, $4/5$ splits improved),
which we do not regard as significant.

Larger temperatures appear to improve accuracy markedly: $\tau = 0.5$
attains RMSE $3.11$ and the lowest NLL of any configuration tested. These
runs are, however, degenerate. The mean variational standard deviation
collapses from $2.8 \times 10^{-1}$ to $1.5 \times 10^{-6}$ and the mean
inclusion probability rises to $0.97$, i.e.\ the posterior contracts to a
point mass over an essentially dense network. Both RMSE and NLL improve
monotonically as this collapse proceeds, so neither predictive metric
identifies the failure; only the posterior scale and the inclusion
probabilities do. We therefore exclude $\tau \geq 0.3$, and note more
generally that predictive accuracy alone is not a safe criterion for
selecting this hyperparameter.

\paragraph{Model-space prior.}
The Beta--Binomial prior on the inclusion indicators, with shape
parameters $(p_a, p_b)$, was varied over seven settings spanning prior
mean inclusion $0.048$ to $0.667$. All runs remained non-degenerate. The
fitted mean inclusion $\alpha$ responds strongly and monotonically,
falling from $0.248$ at our default $p_a = p_b = 1$ to $0.032$ at
$(1, 20)$, i.e.\ sparsity rises from ${\approx}75\%$ to ${\approx}97\%$.
Over that same range test RMSE varies only between $3.43$ and $3.64$,
again within noise. The response is governed by $p_b$ rather than by the
prior mean: $(0.5, 5)$ and $(1, 5)$ share $p_b$ and yield identical
$\alpha = 0.077$, whereas $(0.5, 5)$ and $(1, 10)$ share a prior mean of
$0.091$ and do not; and $(2, 1)$, which raises the prior mean to $0.667$
but leaves $p_b = 1$, reproduces the default $\alpha = 0.248$ exactly.

\paragraph{Consequences.}
Neither knob closes the gap to the sampling-based methods, and we retain
the defaults ($\tau = 0.2$, $p_a = p_b = 1$) throughout. The LBBNN
results reported here should therefore be read as reflecting the
limitations of the variational approximation itself rather than
insufficient tuning.

The sparsity comparison warrants a further caveat. LBBNN's inclusion
level is largely set by the prior and is nearly insensitive to the data:
across \textsc{boston}, \textsc{yacht}, \textsc{energy} and
\textsc{concrete} the fitted $\alpha$ lies in $[0.235, 0.250]$, a spread
of roughly $0.01$, and on a single dataset it can be moved across a
$22$-point range in sparsity at no measurable cost in predictive
accuracy. The sticky samplers, by contrast, recover sparsity levels
between $37\%$ and $62\%$ on the same four datasets, ordered by dataset.
The relevant distinction is thus not which method reports the larger
sparsity figure, but that LBBNN's is specified in advance whereas the
sticky samplers' is inferred.
